\begingroup
\makeatletter\tagsleft@true\let\veqno\@@leqno\makeatother
\ArithmeticHead{343}{R}{{\scriptsize ZUR THEORIE DER REIHE $F(\alpha,\beta,\gamma,x)$.\quad III.}}
\NumberedItem{4}
Eine der merkwürdigsten Relationen finden wir auf folgende Art: Man setze
\[P=P'(1-x)^\mu,\]
\noindent so wird die Differentialgleichung
\[
\begin{aligned}
&\left(\alpha\beta+\frac{\mu(\gamma-(\alpha+\beta+\mu)x)}{1-x}\right)P'\\
&\quad -(\gamma-(\alpha+\beta+2\mu+1)x)\frac{dP'}{dx}-(x-xx)\frac{d^2P'}{dx^2}=0.
\end{aligned}
\]
\noindent Man setze $\mu=\gamma-\alpha-\beta$, so wird diese Gleichung
\[
\begin{aligned}
&(\gamma-\alpha)(\gamma-\beta)P'-(\gamma-(2\gamma-\alpha-\beta+1)x)\frac{dP'}{dx}\\
&\quad -(x-xx)\frac{d^2P'}{dx^2}=0,
\end{aligned}
\]
\noindent woraus leicht folgt
\[P'=F(\gamma-\alpha,\gamma-\beta,\gamma),\]
\noindent also
\begin{equation}
F(\alpha,\beta,\gamma)=(1-x)^{-\alpha-\beta+\gamma}F(\gamma-\alpha,\gamma-\beta,\gamma).
\tag*{(32)}
\end{equation}
\clearpage\SourceLeaf{351}\ArithmeticHead{344}{L}{ANALYSIS.\quad NACHLASS.}
Zur Verwandlung unserer Functionen oder ihrer Verhältnisse in continuirliche Brüche sind folgende Relationen brauchbar:
\begin{equation}
\begin{aligned}
&F(\alpha,\beta,\gamma+1)-F(\alpha,\beta,\gamma)\\
&\quad=-\frac{\alpha\beta x}{\gamma{.}\gamma+1}F(\alpha+1,\beta+1,\gamma+2),
\end{aligned}\tag*{(33)}
\end{equation}
\begin{equation}
F(\alpha,\beta+1,\gamma)-F(\alpha,\beta,\gamma)=\frac{\alpha x}{\gamma}F(\alpha+1,\beta+1,\gamma+1),
\tag*{(34)}
\end{equation}
\begin{equation}
F(\alpha+1,\beta,\gamma)-F(\alpha,\beta,\gamma)=\frac{\beta x}{\gamma}F(\alpha+1,\beta+1,\gamma+1),\text{[*)]}
\tag*{(35)}
\end{equation}
\noindent {\small [S. 41]}
\begin{equation}
\begin{aligned}
&F(\alpha,\beta+1,\gamma+1)-F(\alpha,\beta,\gamma)\\
&\quad=\frac{\alpha(\gamma-\beta)x}{\gamma{.}\gamma+1}F(\alpha+1,\beta+1,\gamma+2),
\end{aligned}\tag*{(36)}
\end{equation}
\begin{equation}
\begin{aligned}
&F(\alpha+1,\beta,\gamma+1)-F(\alpha,\beta,\gamma)\\
&\quad=\frac{\beta(\gamma-\alpha)x}{\gamma{.}\gamma+1}F(\alpha+1,\beta+1,\gamma+2),
\end{aligned}\tag*{(37)}
\end{equation}
\begin{equation}
\begin{aligned}
&F(\alpha-1,\beta+1,\gamma)-F(\alpha+1,\beta-1,\gamma)\\
&\quad=\frac{(\alpha-\beta)x}{\gamma}F(\alpha+1,\beta+1,\gamma+1).
\end{aligned}\tag*{(38)}
\end{equation}

Aus der Verbindung von (36) und (37) folgt
\begin{equation}
\frac{F(\alpha,\beta+1,\gamma+1)}{F(\alpha,\beta,\gamma)}
=\cfrac[l]{1}{1-\cfrac[l]{\dfrac{\alpha(\gamma-\beta)x}{\gamma{.}\gamma+1}}
 {1-\cfrac[l]{\dfrac{(\beta+1)(\gamma+1-\alpha)x}{\gamma+1{.}\gamma+2}}
 {1-\cfrac[l]{\dfrac{(\alpha+1)(\gamma+1-\beta)x}{\gamma+2{.}\gamma+3}}
 {1-\cfrac[l]{\dfrac{(\beta+2)(\gamma+2-\alpha)x}{\gamma+3{.}\gamma+4}}
 {1-\mathrm{etc.}}}}}}
\tag*{(39)}
\end{equation}
\begin{equation}
\frac{F(\alpha+1,\beta,\gamma)}{F(\alpha,\beta,\gamma)}
=\cfrac[l]{1}{1-\cfrac[l]{\dfrac{\beta x}{\gamma}}
 {1-\cfrac[l]{\dfrac{(\alpha+1)(\gamma-\beta)x}{\gamma{.}\gamma+1}}
 {1-\cfrac[l]{\dfrac{(\beta+1)(\gamma-\alpha)x}{\gamma+1{.}\gamma+2}}
 {1-\underset{\mathrm{etc.}}{\dfrac{(\alpha+2)(\gamma+1-\beta)x}{\gamma+2{.}\gamma+3}}}}}}
\tag*{(40)}
\end{equation}
\noindent Also für $\alpha=0$
\par\medskip\hrule\smallskip
{\small\noindent [*) Zwischen die Gleichungen (35) und (36) hat \textsc{Gauss} später die Gleichung hingeschrieben:]
\[
\begin{aligned}
&F(\alpha-1,\beta+1,\gamma)-F(\alpha,\beta,\gamma)\\
&\quad=\frac{(\alpha-\beta-1)x}{\gamma}F(\alpha,\beta+1,\gamma+1).
\end{aligned}
\]
}
\clearpage\SourceLeaf{352}\ArithmeticHead{345}{R}{{\scriptsize ZUR THEORIE DER REIHE $F(\alpha,\beta,\gamma,x)$.\quad III.}}
\begin{equation}
\begin{aligned}
&1+\frac{\beta}{\gamma}x+\frac{\beta{.}\beta+1}{\gamma{.}\gamma+1}xx
 +\frac{\beta{.}\beta+1{.}\beta+2}{\gamma{.}\gamma+1{.}\gamma+2}x^3+\mathrm{etc.}\\[1.2ex]
&\quad=\cfrac[l]{1}{1-\cfrac[l]{\dfrac{\beta}{\gamma}x}
 {1-\cfrac[l]{\dfrac{2(\gamma-\beta)}{\gamma{.}\gamma+1}x}
 {1-\cfrac[l]{\dfrac{(\beta+1)\gamma}{\gamma+1{.}\gamma+2}x}
 {1-\dfrac{3(\gamma+1-\beta)}{\gamma+2{.}\gamma+3}x\ \mathrm{etc.}}}}}
\end{aligned}\tag*{(41)}
\end{equation}
\clearpage
\endgroup
